\section{From weighted traces to crystalline lattices}\label{newgal:section}

Fix a CM field $E$ with maximal real subfield $E^+$, a prime $\ell$ and an
abstract field isomorphism $\iota:k=\overline{\mathbb Q}_\ell\simeq\mathbb C$.
All limits are taken on the $\ell$-adic side. We describe how the representation
of Section~\ref{tr:section} inherits a local condition at a place $w\mid\ell$.
The approximants are finite linear combinations of global characters. Their
coefficients can have arbitrarily large denominators. The local condition concerns
the restriction of the actual global representation, including any extensions in it.

\subsection{The geometric input and the rank of the seed}

In the odd-rank step there is a $d>0$ with $\sqrt{-d}\in E$ and $-d$ a square in
$\mathbb Q_\ell$. The degree $[E^+:\mathbb Q]$ is even, and each nonsplit finite
place of $E/E^+$ is unramified and has residue characteristic different from $\ell$.
Let $a$ be a positive odd integer. A usable place away from $\ell$ is a split
place $w$ whose rational residue characteristic $q$ satisfies
$-d\in(\mathbb Q_q^\times)^2$. Two places are apart when their restrictions
to $E^+$ differ. The auxiliary places $w_0,v'\mid\ell$ are apart. The definite
unitary automorphic representation $U$ is unramified at nonsplit places outside
a given finite set $T$ of nonsplit places. A seed has at $w_0$ a nontrivial
character of a cyclic $\ell$-power quotient of a torsion-free compact open
subgroup, chosen so that every irreducible smooth representation containing
that character is cuspidal. At $v'$ the seed is spherical and has a lacunary
set of $a$ jumps at one embedding. Lacunarity means that two disjoint nonempty
subsets have different arithmetic means.
\src{src/Development/AbelianVarietyTheory/Definitions/GaloisGeometry/LatticeDataThicknessBound.lean}{16}{Field and level assumptions},
\src{src/Development/AbelianVarietyTheory/Definitions/Divisors/GenOf.lean}{14}{local type}, and
\src{src/Development/AbelianVarietyTheory/Definitions/GaloisGeometry/DivisionOfDeterminantCharacter.lean}{15}{lacunarity}.

The geometric theorem has an intermediate output of variable dimension.
One takes a Hecke eigenspace in a cohomology tower, passes to a semisimplification
and twists by the Galois character attached to an algebraic Hecke character.
The resulting representation $V$ is continuous and semisimple. On the usable
spherical primes outside a finite set its geometric Frobenius is annihilated
by the normalized Hecke polynomial of degree $a$. At an auxiliary embedding
at least $a$ lacunary jumps occur. One can request additional geometric
information at one place by changing the tower or the Hecke eigenspace.
The separate requests yield a crystalline $V$ containing the prescribed jumps,
or a crystalline $V$ whose Frobenius polynomial at a degree $-1$ Weil element
divides a power of the predicted polynomial. The geometric statement requires
the specified places to restrict compatibly to $\mathbb Q(\sqrt{-d})$.
Above $\ell$ the two possibilities are exchanged by CM conjugation. Replacing
an auxiliary place by its conjugate preserves apartness at $E^+$, and the
argument uses conjugation of the local type, the spherical condition and
the lacunary condition to make these choices.
\src{src/Development/AutomorphicArithmetic/Foundations/GaloisRepresentations/IntegralStructures/GeometricFlagBases/HodgeTrAt.lean}{79}{Separate geometric outputs}
and \src{src/Development/AutomorphicArithmetic/Foundations/GaloisRepresentations/IntegralStructures/ZetaCharacterDivisibility/SeedFinOddOf.lean}{46}{conjugate-place choice}.

The rank-extraction input takes a continuous semisimple $k$-representation
$V$ of $G_E$ with a lacunary set of $a$ jumps at one labelled embedding
as above. Let $\psi:G_E\to k^\times$ be a continuous character. On a
dense set of $g$ suppose there is a monic polynomial $Q_g$ of degree $a$
annihilating $g$ on $V$ and satisfying $(-1)^a Q_g(0)=\psi(g)$.
The source supplies a continuous semisimple $W$ of dimension $a$ with
equivariant injections
\begin{equation}\label{newgal:injections}
                 W\hookrightarrow V\hookrightarrow W^{\oplus m}
\end{equation}
for some $m$, such that the elements acting with distinct eigenvalues on $W$
form a dense open set. The determinant character in the application comes
from the central character and the Satake normalization. Proposition~\ref{famgal:rank} below proves the rank extraction from the
Hodge cocharacter. The existence of the initial geometric witness is an
arithmetic input here.
\src{src/Development/AutomorphicArithmetic/Foundations/GaloisRepresentations/IntegralStructures/ZetaCharacterDivisibility/PhiOfDivOf.lean}{122}{Rank-extraction hypothesis and conclusion}
and \src{src/Development/AutomorphicArithmetic/Foundations/GaloisRepresentations/IntegralStructures/ZetaCharacterDivisibility/DivisionOneGeometricRepresentation.lean}{52}{division and cocharacter theorems}.

\begin{proposition}\label{newgal:seed}
These inputs identify the same $W$ for all geometric witnesses for a fixed
seed eigensystem. If the witness at $w\mid\ell$ is crystalline and contains
the $a$ prescribed distinct jumps at a labelled embedding, $W$ is crystalline
there and has precisely those jumps. If the crystalline Frobenius polynomial
of the witness divides a power of the predicted polynomial, the predicted
polynomial is the Frobenius polynomial of $W$ whenever the Satake multiset
satisfies the multiplicative genericity condition below.
\end{proposition}
\begin{proof}
At a good geometric Frobenius in the open set of distinct eigenvalues,
the degree $a$ Hecke polynomial annihilates its action on $W\subset V$.
The minimal and characteristic polynomials both have degree $a$, so both
equal the monic Hecke polynomial. For two geometric witnesses, intersect
their dense open sets and discard the union of their finite exceptional
sets. Usable geometric Frobenius elements remain dense there by the density
argument of Section~\ref{tr:section}. Equality of traces extends by continuity
to $G_E$, and semisimplicity in characteristic zero identifies the two $W$'s.
The same argument applied to the constant coefficient gives $\det W=\psi$.
Thus the $W$ already selected admits \eqref{newgal:injections} with each
new witness, up to isomorphism.

Since the global representations are semisimple, both injections split
equivariantly. The splittings persist after restriction to a decomposition
group. Crystalline periods commute with finite direct sums and direct
summands, as do the labelled graded pieces of the de Rham filtration.
Hence $W$ is crystalline when $V$ is, and every jump of $V$ occurs in
$W^{\oplus m}$, hence in $W$. At a fixed embedding the graded dimensions
of $W$ sum to $a$. The occurrence of $a$ distinct prescribed integers forces
multiplicity one at each of them.

Write $S=q_w^{(a-1)/2}B_w$. Multiplicative genericity says that a multiset
of size $a$ supported on $S$ and having product $\prod_{s\in S}s$ equals
$S$, with its multiplicities. Scaling all entries by the same nonzero
scalar preserves this condition. The Frobenius polynomial on $W$ divides
that on $V$ by the splitting. Its roots are supported on $S$, it has
degree $a$, and its root product agrees with that of $S$ by the local
compatibility of the rank-one determinant character. This proves equality
with $\prod_{s\in S}(X-s)$. At a crystalline place the associated Weil
action is unramified, so the characteristic polynomials for all Weil
elements are determined by powers of the degree $-1$ element. Away from
$\ell$ one uses the annihilating polynomial and the same product test.
\src{src/Development/AutomorphicArithmetic/Foundations/GaloisRepresentations/IntegralStructures/ZetaCharacterDivisibility/PhiOfDivOf.lean}{166}{Comparison across geometric realizations}
and \src{src/Development/AutomorphicArithmetic/Foundations/GaloisRepresentations/IntegralStructures/ZetaCharacterDivisibility/PRootPin.lean}{18}{root multiplicities}.
\end{proof}

\subsection{Extracting the rank from a lacunary Hodge type}

\begin{proposition}[Rank extraction in every component]\label{famgal:rank}
Let \(\Gamma\) be profinite and \(V\) a continuous semisimple representation
over \(k=\overline{\mathbf Q}_\ell\). Suppose that on a dense subset the
operator on \(V\) is annihilated by a monic polynomial of degree \(a\ge1\).
Suppose that the algebraic monodromy group of \(V\), after scalar extension
to \(\mathbf C_\ell\), contains a cocharacter for which \(a\) distinct
weights occur. Assume that no two disjoint nonempty subsets of these
\(a\) integers have the same arithmetic mean. There is a semisimple
\(k\)-representation \(W\) of dimension \(a\) with
\(W\hookrightarrow V\hookrightarrow W^{\oplus h}\) for some \(h\).
The elements of \(\Gamma\) having simple spectrum on \(W\) form a dense
open subset.
\end{proposition}
\begin{proof}
Write \(G\) for the Zariski closure of the image and \(H=G^\circ\).
Every element \(A\) of \(G\) has minimal polynomial of degree at most
\(a\). Indeed, linear dependence of \(1,A,\ldots,A^a\) is the
vanishing of minors, and it holds first on the given dense set.
These equations persist over \(\mathbf C_\ell\). The representation of
\(G\) is semisimple since its invariant subspaces are exactly those of
\(\Gamma\). The unipotent radical of \(H\) is normal in \(G\), and its
fixed space in every irreducible \(G\)-summand is nonzero by triangularization.
It thus acts trivially on the faithful representation \(V\), so \(H\)
is reductive.

Choose a maximal torus \(T\) and a Borel subgroup containing it over \(k\).
By conjugacy of maximal tori, we may evaluate the given cocharacter on
the \(T\)-weights after extending scalars. At a general point of \(T\)
all distinct weights give distinct eigenvalues, since their pairwise
quotients are nontrivial characters. There are at most \(a\) weights
and the cocharacter separates at least \(a\). Thus there are exactly
\(a\), all separated, and their cocharacter values are the given set.

In an irreducible \(H\)-module all weights differ by elements of the
root lattice. Averaging over the Weyl group annihilates each root,
since its reflection negates it. The average of each Weyl orbit in
this module is consequently the same rational character. Two different
orbits would give disjoint nonempty subsets of the cocharacter weights
with equal means. There is therefore just one orbit. Its highest weight
has multiplicity one, and Weyl invariance of multiplicities shows that
every weight in the orbit has multiplicity one. If two irreducibles
share a weight they share the orbit and its dominant highest weight,
so they are isomorphic. We use here the ordinary highest-weight theorem
for connected reductive groups in characteristic zero.

We check the action of the components on the multiplicity spaces.
Let \(n\in G(k)\) normalize \(T\) and the chosen Borel. It permutes
the \(a\) weights. For a cycle \(C\) of length \(d_C\), put
\(w_C=\prod_{\alpha\in C}\alpha\). On one weight space \(V_\alpha\)
in the cycle, \((tn)^{d_C}\) acts as \(w_C(t)n^{d_C}\), \(t\in T(k)\).
The spectrum of \(tn\) on the cycle contains all \(d_C\)-th roots of
\(w_C(t)c\) for each eigenvalue \(c\) of \(n^{d_C}|_{V_\alpha}\).
Indeed, propagate an eigenvector around the cycle, dividing by the
chosen root at each step. The resulting components close up on return.
Two different cycles \(C,D\) have disjoint spectra for general \(t\).
An overlap would impose
\(w_C(t)^{d_D}/w_D(t)^{d_C}=c_D^{d_C}/c_C^{d_D}\) for a pair of eigenvalues.
The character on the left is nontrivial by the distinct arithmetic
means on the cocharacter. There are only finitely many pairs.
Thus \(tn\) has at least \(\sum_C d_C=a\) distinct eigenvalues.
Its minimal polynomial has degree at most \(a\), so it is diagonalizable
with exactly this many eigenvalues. Each \(n^{d_C}|_{V_\alpha}\)
is diagonalizable and has only one eigenvalue, hence is scalar.
In particular \(n\) acts by a scalar on any weight space that it fixes.

In each orbit of \(G/H\) on the \(H\)-isotypic spaces choose one
representative and let \(\alpha\) be its highest weight.
The stabilizer of this isotypic space preserves the \(H\)-submodule
generated by any vector of \(V_\alpha\). Indeed, write a stabilizing
element as \(h n\) with \(h\in H\) and \(n\) normalizing the Borel
and torus, using conjugacy of Borel--torus pairs in \(H\).
The element \(n\) fixes \(\alpha\) and acts by a common scalar on
\(V_\alpha\). An \(H\)-intertwiner taking one highest vector to another
also commutes with the stabilizer: on \(h'v\) this follows by writing
\(n h'v=(n h'n^{-1})nv\). A basis of \(V_\alpha\) therefore gives
isomorphic stabilizer modules whose sum is the chosen isotypic space.
Transport each of these modules to the other isotypic spaces in its
orbit. Choice of coset representatives has no effect on the transported
subspaces, and the same intertwiners transport consistently.
Their orbit sums give isomorphic \(G\)-modules. In each orbit retain
one such sum and take their direct sum for \(W\). Every \(T\)-weight
now occurs exactly once. This gives \(\dim W=a\) and the two
injections, taking \(h\) to be the maximum multiplicity.
As a summand of \(V\), \(W\) is semisimple and continuous.

Every component has a representative \(n\) normalizing the Borel and
torus. The cycle calculation on its one-dimensional weight spaces
gives simple spectrum on \(W\) for general \(tn\). The nonvanishing
of the discriminant is thus a nonempty Zariski open set in every
component. The preimage of \(H\) in \(\Gamma\) is open and closed.
Its image is Zariski dense in \(H\), as is the image of each finite-index
subgroup of it, since an irreducible variety cannot be a finite union
of proper closed subsets. Each open subset of \(\Gamma\) contains a
coset of such an open subgroup. Its image meets the discriminant open
in its component. Continuity gives openness on \(\Gamma\) as well.
\end{proof}

For the seed above, use the labelled Hodge--Tate decomposition of the
geometric witness at its distinguished embedding. Its grading defines
a cocharacter in the local algebraic monodromy group, hence in \(G\).
To see the membership, the Hodge--Tate decomposition respects tensor
products, duals, morphisms and subrepresentations of Hodge--Tate
representations. The grading therefore preserves each stable line in
tensor constructions, including direct sums of such constructions.
The stabilizer description of an algebraic subgroup of a general linear
group (Chevalley's theorem) puts the grading in the scalar extension of
that subgroup. With finite coefficients, take the component of
\(\mathbf C_\ell\otimes_{\mathbf Q_\ell}L\) at the embedding giving the
specified label, extending the local embedding to \(\mathbf C_\ell\)
accordingly. This is the labelled cocharacter construction of
Dospinescu, Pa\v{s}k\=unas and Schraen~\cite{senCharacters}.
Changing the sign convention for all weights preserves lacunarity.
Proposition~\ref{famgal:rank} now supplies \eqref{newgal:injections}
for every such witness. The determinant identification is then the
dense Frobenius calculation in Proposition~\ref{newgal:seed}.
The construction still needs the geometric witness and the simultaneous
Hecke congruences at the claimed weights and places.
\src{src/Development/AutomorphicArithmetic/Foundations/GaloisRepresentations/IntegralStructures/ZetaCharacterDivisibility/HtCocharOf.lean}{20}{Hodge cocharacter in the closure},
\src{src/Development/EtaleCohomology/DivisorCycles/GoalHolds.lean}{20}{algebraic rank-extraction input}, and
\src{src/Development/AutomorphicArithmetic/Foundations/GaloisRepresentations/IntegralStructures/ZetaCharacterDivisibility/PDivSet.lean}{16}{application to the Galois seed}.


\subsection{Integral quotients of the regular image}

Let $\Gamma$ be a compact Hausdorff group, $r:\Gamma\to\mathrm{GL}_a(k)$
a continuous semisimple representation and $(r_i)$ a family of continuous
$a$-dimensional representations over $k$. Give $k$ the norm $|\ell|=\ell^{-1}$.
The approximation needed for crystalline closure is
\begin{equation}\label{newgal:lin}
 \inf_{I,c}\ \sup_{g\in\Gamma}
 \left|\operatorname{tr}r(g)-\sum_{i\in I}c_i\operatorname{tr}r_i(g)\right|=0 ,
\end{equation}
where $I$ is a finite subfamily and $c_i\in k$. Fix a finite extension
$K/\mathbb Q_\ell$ and a continuous homomorphism $G_K\to\Gamma$. In the
application $K=E_w$ and this homomorphism identifies $G_K$ with a decomposition
subgroup. Write $V_i,V$ for the restricted modules.

\begin{lemma}[Finite coefficient fields]\label{newgal:coeff}
Each continuous representation of $\Gamma$ over $k$ can be written, in its
given basis, over a finite extension of $\mathbb Q_\ell$.
\end{lemma}
\begin{proof}
There are countably many finite subextensions of $k/\mathbb Q_\ell$.
For each degree there are finitely many local extensions up to isomorphism
by the local-field finiteness theorem, and each has finitely many embeddings.
For each such $L$ the compact image group meets $\mathrm{GL}_a(L)$ in a
closed subgroup, since $L$ is complete. These subgroups cover the image.
Baire's theorem gives one with nonempty interior. It is open and of finite
index. Adjoin the entries of finitely many coset representatives to $L$.
All matrices are defined over the resulting finite extension.
\end{proof}

Choose a field of definition $L$ for $r$. The sum of the translates of any
$\mathcal O_L$-lattice is bounded by compactness and contains the original
lattice. It is therefore a stable lattice. Choose its basis, and do the same
for each $r_i$ over a field $L_i$. Define
\[
 R= r(\mathbb Z_\ell[\Gamma])\subset M_a(\mathcal O_L),\qquad
 R_I=(r_i)_{i\in I}(\mathbb Z_\ell[\Gamma])
       \subset\prod_{i\in I}M_a(\mathcal O_{L_i}).
\]
These are finite free $\mathbb Z_\ell$-modules, stable under left translation
by $\Gamma$. We use the ordinary group ring with finite supports. Both image
modules are already closed, since finite free $\mathbb Z_\ell$-modules are compact.

\begin{lemma}[Lattice quotients from translated traces]\label{newgal:quot}
For every positive integer $n$ there is a finite $I$ and a
$\Gamma$-equivariant surjection $R_I\twoheadrightarrow R/\ell^nR$,
given by $((r_i(f))_i)\mapsto r(f)\bmod\ell^nR$.
\end{lemma}
\begin{proof}
In the $k$-algebra spanned by $r(\Gamma)$ the matrix trace pairing is
nondegenerate by Proposition~\ref{tr:realize}. Choose finitely many $h_j$
so that $r(h_j)$ span this algebra over $k$. On the finite-dimensional
$\mathbb Q_\ell$-space $R[1/\ell]$ the map
\[
   x\longmapsto(\operatorname{tr}(xr(h_j)))_j\in L^b
\]
is therefore injective. The maximum of the norms of its coordinates is
a norm. Comparison with the lattice norm having unit ball $R$ gives an
integer $c\ge0$, depending only on $r$, such that
\[
 \max_j|\operatorname{tr}(xr(h_j))|\le\ell^{-n-c}
       \quad\Longrightarrow\quad x\in\ell^n R.
\]
Use \eqref{newgal:lin} with error at most $\ell^{-n-c}$.
If $f=\sum b_g[g]$ maps to zero in $R_I$, then
$\sum_g b_g\operatorname{tr} r_i(gh_j)=0$ for every $i,j$.
The ultrametric inequality and $|b_g|\le1$ give
$|\operatorname{tr}(r(f)r(h_j))|\le\ell^{-n-c}$ for all $j$.
This proves well-definedness. The formula gives both surjectivity and
equivariance. The kernel contains $\ell^nR_I$, and so is a full stable
lattice in $R_I[1/\ell]$. Global semisimplicity was used before restriction
to $G_K$.
\end{proof}

\begin{theorem}[Crystalline closure from weighted traces]\label{newgal:liu}
Suppose all $V_i$ are crystalline over $K$ and all their labelled filtration
jumps belong to a common finite interval $[u,v]$. Under \eqref{newgal:lin},
$V$ is crystalline over $K$ with jumps in the same interval.
The assertion also holds throughout upon replacing crystalline by semistable.
\end{theorem}
\begin{proof}
Use the usual Fontaine period rings for $K$. With left multiplication on
matrices, the ambient $\mathbb Q_\ell[G_K]$-module of $R_I[1/\ell]$
is a direct sum of $a$ copies of each $\operatorname{Res}_{L_i/\mathbb Q_\ell}V_i$,
using models over the finite extensions $L_i/\mathbb Q_\ell$. Forgetting coefficients collects the jumps
at the different embeddings. In particular the common interval remains valid.
Crystalline (respectively semistable) representations with jumps in that
interval are stable under direct sums and subquotients by exactness and
strictness of the filtered period functor on the admissible category.
Consequently $R_I[1/\ell]$ has this property. Lemma~\ref{newgal:quot}
exhibits $R/\ell^nR$ as a quotient of two full $G_K$-stable $\mathbb Z_\ell$-lattices
inside that single $\mathbb Q_\ell$-representation for each $n$. The ambient
dimension may grow with $n$.

Apply the bounded torsion theorem of T. Liu to $T=R$. In Liu's
contravariant convention the weights are those of the filtered module
of the dual, so our interval becomes $[-v,-u]$. A common integral Tate
twist translates it to $[0,v-u]$. Twisting the quotient sequences uses
the same rank-one $\mathbb Z_\ell$-lattice on both sides. Liu's result
yields the indicated property for $R[1/\ell]$ after undoing the twist.
It allows an arbitrary finite extension $K/\mathbb Q_\ell$, arbitrary
positive height, and $\ell=2$ with the author's correction. For $u=v$ apply
the positive-height statement to both $[u-1,u]$ and $[u,u+1]$ and intersect.
T. Liu,
\emph{Torsion $p$-adic Galois representations and a conjecture of Fontaine},
proved Conjecture 1.0.1, pp.~633--634~\cite{liuTorsion} and
Liu,
\emph{A note on lattices in semi-stable representations}, errata for the 2007 paper~\cite{liuLattices}.

For any $v_0\in L^a$, evaluation $x\mapsto xv_0$ on $R[1/\ell]$ is
$G_K$-equivariant. Taking $v_0$ through a $\mathbb Q_\ell$-basis of $L^a$
yields a surjection from a finite sum of copies of $R[1/\ell]$ onto $L^a$,
since $1\in R$. The quotient is the full restriction of $r$.
The $L$-action commutes with the period comparison map, giving the
coefficient version of the assertion.
\end{proof}

For the automorphic limit take the subfamily indexed by the singleton
$\{w\}$, where $w\mid\ell$ is spherical and apart from $w_0,v'$.
At every embedding $\tau:E\to k$ inducing $w$ the weight at $\iota\tau$
agrees throughout the subfamily. In the zero-based indexing of the highest
weight $\lambda_{\iota\tau}$ the prescribed jumps are
\begin{equation}\label{newgal:jumps}
             J_\tau=\{j-\lambda_{\iota\tau,j}:0\le j<a\}.
\end{equation}
Indeed the $j$th infinitesimal-character entry equals
$\lambda_{\iota\tau,j}+(a-1)/2-j$. Taking the minimum and maximum of
\eqref{newgal:jumps} over $j$ and the finitely many $\tau$ over $w$
gives the common interval. The singleton subfamily retains polynomial
approximation at every precision, hence \eqref{newgal:lin} on the whole
global group. Theorem~\ref{newgal:liu} therefore proves crystallinity
of the limit at this $w$.
\src{src/Development/AutomorphicArithmetic/Foundations/GaloisRepresentations/TraceLimits/LimFinOddOf.lean}{166}{Singleton restriction}
and \src{src/Development/AbelianVarietyTheory/Proofs/SatakeParameters/BlowUpFibreSemilocal.lean}{29}{jump convention}.

\subsection{Comparison of exact local invariants}

For exact jumps and Frobenius the source keeps more information than the
surjection of Lemma~\ref{newgal:quot}. Let $b$ be a $\mathbb Z_\ell$-basis
of $R$, let $e$ be the column of coordinates of $1$, and let $t$ be the
row of traces of the $b_j$. For precision $n$ choose the approximation of
Lemma~\ref{newgal:quot}. Enlarge its support $I_n$ to contain all previous
supports, assigning coefficient zero to the added members. Set
\[
 D_n=\prod_{i\in I_n}M_a(\mathcal O_{L_i}),\qquad
 B_n=R_{I_n}+\ell^{M_n}D_n .
\]
Choose increasing $M_n$ with $\ell^{-M_n}\max(1,\max_i|c_i|)
\le\ell^{-n-c}$. If $(r_i(f))_i\in\ell^{M_n}D_n$, then the
translated traces of $r(f)$ against all $r(h_j)$ have norm at most
$\ell^{-n-c}$ by \eqref{newgal:lin} and integrality of all $r_i(h_j)$.
The map of Lemma~\ref{newgal:quot} thus extends to $B_n$ by zero on
$\ell^{M_n}D_n$. Choose a basis of $B_n$ and any integral matrix $P_n$
lifting that map. Write $e_n$ for the coordinates of the tuple of identities,
and $t_{ni}$ for the row of traces after projection of the basis to factor $i$.
Then
\begin{equation}\label{newgal:marked}
 |P_ne_n-e|\le\ell^{-n},\qquad
 \left|tP_n-\sum_{i\in I_n}c_i t_{ni}\right|\le\ell^{-n}.
\end{equation}
Here matrix norms are maximum-entry norms. The first bound follows
by evaluating the quotient map at $1$. For the second, write each basis
vector as $(r_i(f))_i+\ell^{M_n}z$. Its chosen lift agrees with $r(f)$ modulo
$\ell^nR$, whose traces have norm at most $\ell^{-n}$. The approximation
applied to $f$, and the choice of $M_n$ applied to $z$, give the bound.
The coordinate projection $B_{n+1}\to B_n$ intertwines the actions exactly
and commutes with the maps to $R/\ell^n R$. Both claims follow on group
elements and on $\ell^{M_{n+1}}D_{n+1}$. These constructions work after
restriction to $G_K$ and are compatible with any additional scalar
approximated by the same coefficients.
\src{src/Development/AbelianVarietyTheory/Proofs/InertiaEigenvalues/CoverExists.lean}{23}{Integral model with trace and unit coordinates}.

\begin{proposition}[Reading a labelled Hodge multiset]\label{newgal:hodge}
Fix an embedding $\tau:K\to k$. Suppose the comparison on the preceding
lattices supplies a non-root of unity $z\in\mathbb C_\ell^\times$ and matrices
$U,U_n$ on $R,B_n$ over $\mathbb C_\ell$ with the following properties.
Their norms are bounded by a common $B\ge1$. For every $N$ there is $n\ge N$
with $|UP_n-P_nU_n|\le B\ell^{-N}$. After the same isometric coefficient
embedding on all trace rows and all $c_i$, for $j\ge0$ one has
\[
 tU^je=\sum_{h\in{\rm HT}_\tau(V)}z^{jh},\qquad
 t_{ni}U_n^je_n=\sum_{h\in J_\tau}z^{jh}\quad(i\in I_n).
\]
Here $\operatorname{HT}_\tau$ denotes the multiset of labelled filtration jumps.
Then $\operatorname{HT}_\tau(V)=J_\tau$.
\end{proposition}
\begin{proof}
Put $S_j=\sum_{h\in J_\tau}z^{jh}$. For fixed $j$, the telescoping formula
for $U^jP_n-P_nU_n^j$ bounds its norm by $B^j\ell^{-N}$ for $j\ge1$.
The maximum-entry norm is submultiplicative in all dimensions here.
In $tU^je$ replace $e$ by $P_ne_n$, commute $P_n$ past $U^j$, then
replace $tP_n$ by $\sum c_it_{ni}$. Equations~\eqref{newgal:marked}
bound the total error by $\max(1,|t|)B^j\ell^{-N}$. Thus the result is
$S_j\sum c_i$ to that precision. At $j=0$ both power sums equal $a$,
so $|1-\sum c_i|=O(\ell^{-N}/|a|)$. The same coefficients serve at
$j$ and at zero. Let $N$ tend to infinity. Newton's identities now
identify the multisets $\{z^h:h\in{\rm HT}_\tau(V)\}$ and
$\{z^h:h\in J_\tau\}$. The powers of $z$ distinguish the integers.
\end{proof}

The period-theoretic hypothesis in Proposition~\ref{newgal:hodge} is
quantitative. In the source the height is bounded by
$H=\max_{\tau,h\in J_\tau}|h|$. For its Breuil--Kisin period
$u\in A_{\rm inf}$ put $s_0=\theta(u^{2H})\ne0$, where
$\theta:A_{\rm inf}\to\mathcal O_{\mathbb C_\ell}$ is Fontaine's map.
The linked frames give integral evaluated change-of-basis matrices
$Y,Y'$ with integral matrices $Y_1,Y'_1$ satisfying
$YY_1=Y'Y'_1=s_0\,1$ in the respective dimensions.
For integral cocycle matrices $C,C'$ the operators satisfy $UY=YC$,
$U_nY'=Y'C'$. A link at level $N$ gives an integral rectangular matrix
$D$ with $|P_nY'-YD|,\ |CD-DC'|\le\ell^{-N}$. Thus
$|U|,|U_n|\le |s_0|^{-1}$. The identity
\[
 (UP_n-P_nU_n)Y'=U(P_nY'-YD)+Y(CD-DC')-(P_nY'-YD)C'
\]
and multiplication by $Y'_1/s_0$ give the claimed bound with
$B=|s_0|^{-2}$, independent of the growing dimensions.
The other period assertion identifies the operator on a Hodge-adapted
basis with a diagonal matrix of powers $z^h$ and transports its powers
through each column projection. The construction of those linked
bounded-height frames and of the adapted comparison, needed at every
precision and for the selected embedding, are cited here.
\src{src/Development/AbelianVarietyTheory/Proofs/InertiaEigenvalues/ExistsDeRhamFrameLimitLatticeBreuilKisinTower.lean}{16}{Linked frames},
\src{src/Development/AbelianVarietyTheory/Proofs/InertiaEigenvalues/ThetaConjBoundedCompatibleBreuilKisinFrames.lean}{117}{uniform estimate}, and
\src{src/Development/AbelianVarietyTheory/Proofs/InertiaEigenvalues/ClosedOfCover.lean}{159}{labelled comparison}.

At a crystalline place put $D_i=\iota\mathrm{WD}(V_i)$ and
$D=\iota\mathrm{WD}(V)$, and set $p_i=\iota^{-1}\operatorname{tr}D_i(g)$
and $p=\iota^{-1}\operatorname{tr}D(g)$ for $g$ in the Weil group of $K$.
Write $B_w$ for the spherical Satake
multiset. At $g^j$ the predicted value is
\[
 \vartheta_j=\iota^{-1}\left(\sum_{x\in q_w^{(a-1)/2}B_w}x^{-j\deg g}\right).
\]
The Hecke approximation at each fixed $j$ uses the same $c_i$ for the global
trace and this additional scalar, to arbitrary $\ell$-adic precision.
One removes a fixed nonzero normalization factor by increasing the required
precision. For the next proposition write $\vartheta$ for the target at $g$.

\begin{proposition}[Comparison of a crystalline Frobenius trace]\label{newgal:phi}
Assume the marked approximations \eqref{newgal:marked} satisfy also
$|\vartheta-\sum_{i\in I_n}c_i p_i|\le\ell^{-n}$. Put
$F=\mathbb Q_\ell^{\mathrm{nr}}\subset k$ and use the unramified embedding
$F\hookrightarrow B_{\rm cris}$ taking $\mathcal O_F$ into $A_{\rm cris}$.
Suppose there is a single $s\in A_{\rm cris}$ invertible in $B_{\rm cris}$,
and matrices $A,A_n$ over $A_{\rm cris}$ on the respective lattices, with
the following two properties. For every $N$ some $n\ge N$ has every entry
of $AP_n-P_nA_n$ in $\ell^NA_{\rm cris}$. For every $F$-linear
$\lambda:k\to F$, acting entrywise on the rows, one has
\[
 \lambda(t)Ae=s\lambda(p),\qquad
 \lambda(t_{ni})A_ne_n=s\lambda(p_i)\quad(i\in I_n)
\]
in $B_{\rm cris}$. Then $p=\vartheta$.
\end{proposition}
\begin{proof}
Suppose $\Delta=p-\vartheta\ne0$. Fix integers $c_1,c_2\ge0$ such that
$\ell^{c_1}t$ is integral and $\ell^{-c_2}\le|\ell^{c_1}\Delta|\le1$.
Choose $n$ as in the first hypothesis for $N$ and write
$\delta=tP_n-\sum_i c_i t_{ni}$ and $\eta=\vartheta-\sum_i c_i p_i$.
Both have norm at most $\ell^{-N}$. There is an $F$-linear $\lambda$
taking integral values on the entries of $\ell^{c_1}t,\ell^{-N}\delta,
\ell^{-N}\eta$ and on $z=\ell^{c_1}\Delta$, with $|\lambda(z)|\ge|z|$.
Indeed their $\mathcal O_F$-span is finite free over the discrete valuation
ring $\mathcal O_F$. All elements of that span are integral in $k$. In any
basis of the span one coordinate of $z$ has absolute value at least $|z|$.
Extend this coordinate functional $F$-linearly to $k$.

Apply the seed identities to $x\mapsto\lambda(c_i x)$, subtract and
insert $P_n$. This gives
\[
 \begin{split}
 s\lambda(\Delta)={}&\lambda(t)\big[(AP_n-P_nA_n)e_n+A(e-P_ne_n)\big]\\
                    &+\lambda(\delta) A_ne_n-s\lambda(\eta).
 \end{split}
\]
Multiplication by $\ell^{c_1}$ puts the right side in $\ell^NA_{\rm cris}$.
Multiplication next by $\ell^{c_2}/\lambda(z)\in\mathcal O_F$ gives
$\ell^{c_2}s\in\ell^NA_{\rm cris}$ for every $N$.
This contradicts $\ell$-adic separation and torsion-freeness of
$A_{\rm cris}$, since $s$ is invertible in $B_{\rm cris}$.
\end{proof}

Here is the integral calculation giving the first Frobenius comparison
from the period frames. Use a target precision \(N\) and perform the
calculation in \(A_{\rm cris}/\ell^N\). For the two Frobenius iterates the
linked period bases give square matrices \(W,Z\) and \(W',Z'\) on the
target and seed-product lattices, respectively. Their scaled inverses
satisfy \(Z^*Z=b1,\ Z'Z'^*=b1\) for the same scalar \(b\). The two
frame congruences read \(PZ'=ZD,\ PW'=WD\), where \(P=P_n\) and \(D\)
have size \(\operatorname{rank}R\) by \(\operatorname{rank}B_n\).
In this quotient,
\[
 bWZ^*P=WZ^*PZ'Z'^*=WZ^*ZDZ'^*=bWDZ'^*=bPW'Z'^*.
\]
Thus the integral operators \(A=bWZ^*,\ A_n=bW'Z'^*\) intertwine to
precision \(N\), without inverting \(P\) or cancelling \(b\).
In the construction, if \(H\) bounds the height, the integer \(k_0\)
for the period denominators is chosen before the level. With
\(\tau_1\) the image of the chosen period in \(A_{\rm cris}\), one has
\[
 b=\varphi^{n_2}(\tau_1^{2H}\ell^{2k_0}),\qquad
 s=\varphi^{n_2}(\tau_1)^{3H}\ell^{4k_0}\varphi^{n_1}(\tau_1)^H.
\]
Both lattice dimensions use these same scalars. To obtain the two frame
congruences modulo \(\ell^N\), first ask for the frame-link precision
\(N'\) required by the period comparison, then choose one level
\(n\ge\max(N',N)\). At that level both iterations compare with the
same \(D\). The unit error, weighted trace-row error and the error in
the additional scalar are already at most \(\ell^{-n}\). In particular
one never multiplies an error of size \(\ell^{-N}\) for an individual
seed by \(c_i\). The potentially unbounded \(\max_i|c_i|\) entered
earlier in the choice of \(M_n\) in \eqref{newgal:marked}.
\src{src/Development/AutomorphicArithmetic/Foundations/GaloisRepresentations/ClosedPhiFunctoriality/OpExHolds.lean}{75}{Common denominator before the level},
\src{src/Development/AutomorphicArithmetic/Foundations/GaloisRepresentations/ClosedPhiFunctoriality/OpExHolds.lean}{86}{one level for both iterates}
and \src{src/Development/AutomorphicArithmetic/Foundations/GaloisRepresentations/ClosedPhiFunctoriality/LinkOp.lean}{8}{the integral rectangular identity}.

One can also see exactly where the unit and trace marks enter the
second Frobenius comparison. For either lattice let \(b_i\) be its
matrix readings (on a product lattice, project to one seed factor);
thus \(\sum_i e_i b_i=1_a\), \(t_i=\operatorname{tr}(b_i)\).
Work after extending the Weil module to the common coefficient and
period ring \(k\otimes_F B_{\rm st}\), writing the two inclusions of
scalars implicitly. Write \(U_j=\varphi^{n_j}(U)\) for the invariant
period basis of this lattice, so \(AU_2=sU_1\). Denote the \(h\)th
standard vector in coefficient dimension \(a\) by \(u_h\).
The column comparison for the crystalline periods supplies vectors
\(v_{jh}\) in the Weil module, whose comparison coordinates are the
columns \(\sum_i (U_2)_{ij}b_i u_h\). Their transforms under \(g\)
have coordinates \(\sum_i (U_1)_{ij}b_i u_h\). This is the period
input identifying the sign of the Weil action. Now set \(z=U_2^{-1}e\).
For each \(h\) the vector \(\sum_j z_jv_{jh}\) has comparison
coordinates \(\sum_i e_i b_i u_h=u_h\).
Applying the scalar-extended Weil operator and summing its diagonal
coordinates gives
\[
 p\otimes 1=\sum_i \operatorname{tr}(b_i)\otimes(U_1z)_i,\qquad
 p\otimes s=\sum_i t_i\otimes(Ae)_i .
\]
Applying \(\lambda\otimes1\) for any \(F\)-linear \(\lambda:k\to F\)
gives precisely \(\lambda(t)Ae=s\lambda(p)\); the terms belong to
\(B_{\rm cris}\subset B_{\rm st}\). The same proof for each column
projection of the product lattice gives every seed identity before
weighting it. Only the column comparison and the invariant periods
are used here beyond linear algebra.
\src{src/Development/AutomorphicArithmetic/Foundations/GaloisRepresentations/ClosedPhiFunctoriality/TieRep.lean}{75}{The two systems of period columns},
\src{src/Development/AutomorphicArithmetic/Foundations/GaloisRepresentations/ClosedPhiFunctoriality/TieRep.lean}{94}{The unit frame and trace contraction}
and \src{src/Development/AutomorphicArithmetic/Foundations/GaloisRepresentations/ClosedPhiFunctoriality/ClosedPhiOfClosedTr.lean}{65}{Application to the target and each product projection}.

In the crystalline comparison the two nonnegative Frobenius iterates are
$n_1=\max(0,-\alpha)$ and $n_2=\max(0,\alpha)$, where
$\alpha=[k(w):\mathbb F_\ell]\deg g$ is the signed exponent on the
unramified coefficient field and $k(w)$ denotes the residue field of $K$.
For an invertible matrix $U$ of invariant periods one requires
$A\varphi^{n_2}(U)=s\varphi^{n_1}(U)$, and likewise at each level.
The common $s$ and congruences of these operators are assertions about
the linked integral frames. The second hypothesis of the proposition
identifies their values on the marked unit, paired with the marked trace
rows, with the Weil trace via the period comparison. Using $F$ and its
linear functionals avoids choosing embeddings of all ramified coefficient
fields into $B_{\rm cris}$. Given these two frame assertions, apply
Proposition~\ref{newgal:phi} at the first $a$ positive powers of a
geometric Frobenius. Newton's identities give the complete characteristic
polynomial. Crystallinity gives trivial inertia and zero monodromy,
so this determines the Frobenius-semisimplified Weil--Deligne parameter.
\src{src/Development/AutomorphicArithmetic/Foundations/GaloisRepresentations/ClosedPhiFunctoriality/IsOp2One.lean}{45}{Integral operators and common denominator},
\src{src/Development/AutomorphicArithmetic/Foundations/GaloisRepresentations/ClosedPhiFunctoriality/IsOp2One.lean}{81}{comparison with Weil traces}, and
\src{src/Development/AutomorphicArithmetic/Foundations/GaloisRepresentations/ClosedPhiFunctoriality/ClosedPhiOfClosedTr.lean}{13}{passage through the integral comparison}.

In particular, Theorem~\ref{newgal:liu} supplies an interval for the
weights and the existence of the crystalline period space. It gives no
formula pairing $t$ with a specific operator on that space. A finer
published congruence theorem for odd $\ell$ gives equality of the labelled
Hodge types of two semistable representations of the same rank over one
$\mathcal O_L$ if their matrices are congruent modulo $\ell^n$, with
$n>c(K,H)$ for a common positive height bound $H$ after twisting.
Here $P_n$ can be rectangular and the approximation of $t$ uses a
weighted sum. Applying that theorem to $P_n$ would require an additional
equal-rank congruence. Propositions~\ref{newgal:hodge} and~\ref{newgal:phi}
specify the uniform comparisons used for this more general family. An
independent construction of the linked integral frames with the two
period identifications is still required for this proof in that generality.
T. Liu,
\emph{Filtration associated to torsion semi-stable representations},
Theorems 1.1 and 4.16~\cite{liuFiltration}.

In particular, at \(\ell=5\) the trivial and the unramified quadratic
character have the same single Hodge jump \(0\), while their Frobenius
traces are \(1\) and \(-1\). The local approximation at a spherical
\(w\mid\ell\) therefore uses the extra Hecke coordinates explicitly.
Take \(Y=\{w\}\) for \(w\) apart from the two auxiliary places. The
congruence theorem here holds on the disjoint union of the global
indices and the local integer powers. For a fixed Weil element \(g\)
and power \(j\), choose the coordinate with index \(-j\deg g\).
Its nonzero Hecke normalization is fixed before requesting the
precision. Divide it out of the approximated scalar and request the
correspondingly finer error in the original approximation. Genericity
at this \(w\) in every selected seed identifies its local power sum
with its Weil trace. This gives exactly the simultaneous global and local
test of Proposition~\ref{newgal:phi}.
\src{src/Development/AutomorphicArithmetic/Foundations/GaloisRepresentations/TraceLimits/LimFinOddOf.lean}{82}{Joint coordinate set}
and \src{src/Development/AutomorphicArithmetic/Foundations/GaloisRepresentations/TraceLimits/LimFinOddOf.lean}{185}{normalization and the same coefficients}.

The labelled comparison counts multiplicities. At power zero the
augmentation has error at most the prescribed error by testing the
constant polynomial \(1\) in \eqref{tr:approx}. At each fixed power
the same-level rectangular estimate in Proposition~\ref{newgal:hodge}
then recovers the full power sum, including repeated summands. Newton's
identities require no separation of jumps. For a dominant integral
weight in \eqref{newgal:jumps}, even equal adjacent highest-weight
entries give distinct shifted jumps, since
\((j+1-\lambda_{j+1})-(j-\lambda_j)\ge1\). Neither this observation nor
the multiset argument produces an automorphic seed at a singular
intermediate stage. The input there is the weight and local-type
congruence together with the geometric seed and the theta-lift and
descent steps in their stated ranges. For the odd-rank limit the
singleton construction compares weights only away from \(w_0,v'\).
\src{src/Development/AutomorphicArithmetic/Foundations/GaloisRepresentations/TraceLimits/LimFinOddOf.lean}{166}{Range of the singleton comparison}
and \src{src/Development/AutomorphicArithmetic/Foundations/GaloisRepresentations/IntegralStructures/Development.lean}{27}{Assembly from geometry, congruences and lifting}.

\subsection{The classical polarized cuspidal case}

There is a direct proof of the local assertions for the final classical
$\mathrm{GL}_a$ input which also explains the normalization. Let $\pi$
be a cuspidal automorphic representation of $\mathrm{GL}_a(\mathbb A_E)$
whose archimedean infinitesimal character is that of an irreducible
algebraic representation. Assume a continuous Hecke character $\chi$
of $E^+$ satisfies
\[
 \pi^c\simeq\pi^\vee\otimes(\chi\circ N_{E/E^+}\circ\det),
 \qquad \chi_v(-1)=(-1)^a\quad(v\mid\infty).
\]
Here $c$ is CM conjugation. The invariant bilinear pairing in the
polarization gives precisely this duality when the smooth finite and
archimedean modules are those of $\pi$. Choose the classical local
Langlands correspondence with uniformizers mapping to geometric Frobenius.
\src{src/Development/AbelianVarietyTheory/Definitions/GaloisGeometry/DualIsConjByOne.lean}{53}{Polarization}
and \src{src/Development/AutomorphicArithmetic/Foundations/GaloisRepresentations/TypeConditions/WeightsConjAlt.lean}{85}{cuspidal input}.

\begin{proposition}
For this input\label{newgal:classical} there is a continuous semisimple $r_\pi:G_E\to
\mathrm{GL}_a(k)$ with labelled jumps \eqref{newgal:jumps} at all
$w\mid\ell$. It is de Rham there and crystalline when $\pi_w$ is
spherical. Write
\[
 Y_w=\operatorname{rec}_{E_w}\bigl(\pi_w\otimes|\det|_w^{(1-a)/2}\bigr).
\]
At every $w\nmid\ell$ one has $\iota\mathrm{WD}(r_\pi|_{G_{E_w}})^{F\text{-}\mathrm{ss}}
\simeq Y_w$. At $w\mid\ell$ this equality also holds if the local
representation has an Iwahori-fixed vector or if $Y_w$ has zero monodromy.
Any continuous semisimple $r$ with the same normalized traces on a dense
set of good Frobenius elements inherits these conclusions.
\end{proposition}
\begin{proof}
Apply Barnet-Lamb, Gee, Geraghty and Taylor's Theorem~2.1.1 to $(\pi,\chi)$.
Their weight $\mu$ uses the infinitesimal character of the dual of the
algebraic module of highest weight $\mu$. Thus, with their index $1\le i\le a$,
$\mu_{\iota\tau,i}=-\lambda_{\iota\tau,a-i}$. The theorem gives jumps
$\mu_{\iota\tau,i}+a-i$, precisely \eqref{newgal:jumps}. It gives an
integer $w_\pi$ such that the jumps at $\tau c$ are the complements
$w_\pi-h$ of those at $\tau$. Parts (2) and (4) give the asserted
comparisons away from $\ell$ and at Iwahori level respectively, including
crystallinity at spherical level. Twisting by the indicated power of the
norm multiplies spherical Satake eigenvalues at geometric Frobenius by
$q_w^{(a-1)/2}$.
Barnet-Lamb,
Gee, Geraghty and Taylor, \emph{Potential automorphy and change of weight},
\S2.1, Theorem 2.1.1(2)--(4)~\cite{blggtPotential}.

For $a\ge2$ and $w\mid\ell$, the same authors' monodromy bound gives
the semisimplified Weil representation of $Y_w$. In each class of
irreducible Weil constituents up to unramified twist, the decreasing
list of lengths of the special representations on the Galois side is
dominated in partial sums by that on the automorphic side. If the
automorphic monodromy is zero, every length is one. The bound on the
first partial sum forces every Galois length to be one. Hence its
monodromy also vanishes, proving the claimed equality. For $a=1$
the same equality follows from the local compatibility of the
$\ell$-adic character of an algebraic Hecke character. Finally the
two global traces agree by continuity on $G_E$ and the characteristic-zero
Brauer--Nesbitt theorem identifies $r$ with $r_\pi$.
Barnet-Lamb,
Gee, Geraghty and Taylor, \emph{Local-global compatibility for $l=p$, II},
\S1, Theorems 1.1 and 3.1~\cite{blggtLocal}.
\end{proof}

For the classical correspondence these conclusions cover the ramified
places tested in the final statement. Above $\ell$ its extra local
comparison explicitly assumes zero monodromy on the automorphic side.
Away from $\ell$ the full comparison gives equality of all the tested
inertial polynomials and monodromy conditions. Applying this external
argument to intermediate definite-unitary eigensystems would require
their identification with suitable classical $\mathrm{GL}_a$ data and
equality of the specified Hecke systems. Applying it with an abstract
reciprocity map would also require the tested local parameters of that
map to agree with the classical ones. The weighted trace construction
and integral comparisons above describe the source's route prior to
these identifications.

\subsection{The principal level of an attached classical system}

Let $F$ be an imaginary quadratic field in rank three, or $\mathbb Q$ in
rank two, and let $M\ge1$. Write $\varepsilon_\ell$ for the cyclotomic
character. A character $\chi:G_F\to k^\times$ trivial on
$G_{F(\mu_M)}$ has finite image. For $v\mid p$ write $d=v_v(M)$,
where $v_v$ gives a uniformizer value one. Put
$U_v^0=\mathcal O_v^\times$ and $U_v^d=1+\mathfrak p_v^d$ for $d>0$.
Under local class field theory $\chi_v$ is trivial on $U_v^d$.
Indeed $\operatorname{Gal}(F(\mu_M)/F)$ is a subgroup of
$(\mathbb Z/M)^\times$. Extend its character to this finite abelian
group, using divisibility of $k^\times$. The local character on $F_v$
is then the norm pullback of a character on $\mathbb Q_p^\times$ whose
restriction to $\mathbb Z_p^\times$ factors modulo $p^{v_p(M)}$.
When $d>0$, any element of $U_v^d$ acts by multiplication on the
$\mathbb Z_p$-lattice $\mathcal O_v$ through a matrix congruent to
one modulo $p^{v_p(M)}$, since $\mathfrak p_v^d=M\mathcal O_v$.
Taking its determinant proves the claim. When $d=0$ the character
on $\mathbb Q_p^\times$ is unramified.

\begin{lemma}[Fixed vector criterion]
Let $\pi_v$\label{newgal:fixed} be a generic irreducible smooth representation of
$\mathrm{GL}_n(F_v)$, and write $(D,r_W,N)$ for its
Frobenius-semisimple Weil--Deligne parameter, allowing an unramified
normalizing twist. Suppose its inertia representation is a sum of
restrictions of characters of $W_{F_v}$ of conductor exponent at most
$d$. If $d=0$, assume also $N=0$. Then $\pi_v$ has a nonzero vector
fixed by the full principal congruence group $K_v(d)$, where
$K_v(0)=\mathrm{GL}_n(\mathcal O_v)$.
\end{lemma}
\begin{proof}
Write the generic representation as the irreducible parabolic induction
of essentially square-integrable representations $\sigma_j$ of ranks
$n_j$. Their parameters $D_j$ are the corresponding direct summands of
$D$. For $d>0$, set $t_j=\dim D_j^{I_v}$. The Artin conductor of its
Weil part satisfies $a(r_{W,j})\le (n_j-t_j)d$: finite inertia acts
semisimply, and every nontrivial constituent has the same restriction
to the ramification groups as one of the specified characters. The
additivity of the Artin conductor on this restriction gives the bound.
The Weil--Deligne conductor adds
$t_j-\dim(\ker N_j)^{I_v}\le t_j$. Thus
\[
 a(\sigma_j)=a(r_{W,j},N_j)\le (n_j-t_j)d+t_j\le n_jd .
\]
Here equality of the automorphic and Weil--Deligne conductors follows
from equality of their standard epsilon factors with the same additive
character, and an unramified twist preserves the exponent.
Miyauchi and Yamauchi's Theorem~1.2 characterizes $K_v(d)$-invariants
by exactly these inequalities for every inducing block. If $d=0$, all
the conductors vanish since inertia and monodromy are trivial. The
same theorem then gives a spherical vector.
Miyauchi and Yamauchi,
\emph{A remark on conductor, depth and principal congruence subgroups},
Theorems 1.1--1.2~\cite{miyauchiYamauchi}.
\end{proof}

\begin{proposition}[Transfer from the prescribed Galois level]\label{newgal:level}
Suppose a continuous semisimple $r$ is attached on a cofinite set of
places to a regular algebraic cuspidal $\pi$ over $F$. Assume $\pi$
is conjugate self-dual in the rank-three case, and comes from the
normalized classical weight $2e+1$ newform, $e\ge1$, in the rank-two
case. In rank three use $r$ after the inverse cyclotomic twist occurring
in Lemma~\ref{rk3:purity}. Suppose $r$ satisfies the following level
conditions for one prime $\ell$ and a tuple $m$.
\begin{enumerate}
\item At $v\nmid\ell$ the restriction to inertia is diagonal, with entries
$\chi_i|_{I_v}$, where each $\chi_i$ is trivial on $G_{F(\mu_M)}$.
\item At $v\mid\ell$ there is a $G_{F_v}$-stable full flag with graded
characters $\kappa_i$. On inertia $\kappa_i=\varepsilon_\ell^{m_i}\chi_i$,
with $\chi_i$ as above. At this $v$ either $r$ is potentially crystalline
or $\kappa_i\ne\kappa_j\varepsilon_\ell$ for all $i,j$.
\end{enumerate}
Then $\pi_f^{K_F(M)}\ne0$. At each embedding over $v\mid\ell$ its
attached representation has filtration jumps $\{-m_i\}$. In the two
counting applications $m=(e-1,-1,-e-1)$ and $m=(e,-e)$ respectively.
\end{proposition}
\begin{proof}
Chebotarev and semisimplicity identify $r$ with the representation
of $\pi$; this uses attachment to the classical cuspidal system as
a hypothesis. Every $\pi_v$ is generic by global cuspidality for
$\mathrm{GL}_n$. At $v\nmid\ell$ the prescribed finite inertia action
implies zero monodromy, since it is trivial on an open subgroup of
inertia, and gives precisely the specified inertia characters for the
Weil parameter. Full local--global compatibility away from $\ell$
therefore gives Lemma~\ref{newgal:fixed} at such $v$.

At $v\mid\ell$ put $L=F_v$. The graded characters have the form
$\kappa_i=u_i\chi_i\varepsilon_\ell^{m_i}$, with $u_i$ continuous
and unramified. Its value on geometric Frobenius is an $\ell$-adic unit
by compactness of $G_L$. Each graded line has the labelled jump
$-m_i$ at every embedding of $L$, and the attached representation is
de Rham. Strict exactness of filtered de Rham comparison now gives
the stated jumps. Let $q_v$ be the residue cardinality. The Weil
character of $\varepsilon_\ell$ is the unramified character $\nu$
taking geometric Frobenius to $q_v^{-1}$. On a finite character and on
a continuous unramified character the Weil functor gives restriction
to $W_L$, with coefficients transported by $\iota$.
Exactness of potentially semistable comparison gives the
semisimplified Weil part as $\bigoplus_i\mathrm{WD}(\kappa_i)$.
In particular its inertia is the sum of the $\chi_i|_{I_v}$.

We check zero monodromy when $d=0$. Potential crystallinity gives this
immediately in the first case. Otherwise, a nonzero monodromy matrix
entry between the $j$th and $i$th character spaces would give
$\mathrm{WD}(\kappa_i)=\nu\,\mathrm{WD}(\kappa_j)$ by
$r_W(g)Nr_W(g)^{-1}=\nu(g)N$. The character
$\theta=\kappa_i/(\kappa_j\varepsilon_\ell)$ would then have trivial
Weil character. Write $\theta=u\chi\varepsilon_\ell^s$ as above. On
geometric Frobenius, after returning to $k$, the Weil character takes
the value $u(\mathrm{Frob})\chi(\mathrm{Frob})q_v^{-s}$. Its
$\ell$-adic valuation forces $s=0$. The finite and unramified factors
now give $\theta|_{W_L}=1$, hence $\theta=1$ by density. This contradicts
the second alternative in the flag hypothesis.

Finally the needed full comparison at a possibly ramified
$v\mid\ell$ follows in rank three because every odd rank is
Shin-regular. In rank two the difference of the two dominant weight
entries is $2e-1>0$, as obtained by subtracting the half-sum of
positive roots from the normalized infinitesimal character. This is
the Shin-regular condition in rank two. Apply Theorem~1.1 for the CM
case and Corollary~1.2 for the totally real case in
Barnet-Lamb,
Gee, Geraghty and Taylor, \emph{Local-global compatibility for $l=p$, II}~\cite{blggtLocal}.
Thus Lemma~\ref{newgal:fixed} applies to the actual parameter on the
automorphic side, including at the coefficient prime. Taking the restricted
tensor product of the resulting fixed vectors gives a nonzero vector
at principal level $M$. Use $M=N$ in rank three and
$M=N|\operatorname{disc}F'|$ in rank two, where $F'$ is the imaginary
quadratic field to which the rank-two representation is subsequently
restricted. These are precisely the levels used in the fixed-field
count of Proposition~\ref{rk3:count}.
\end{proof}

The exported rank-three counting argument calls both
\src{src/Development/AutomorphicArithmetic/AutomorphicVarieties/ModularityArgument/ThresholdArguments/DenseThresholdCases/KeepLevel.lean}{95}{the conjugate self-dual CM level theorem}
and \src{src/Development/AutomorphicArithmetic/AutomorphicVarieties/ModularityArgument/ThresholdArguments/DenseThresholdCases/KeepLevel.lean}{83}{the polarized CM level theorem},
as recorded at the
\src{src/Development/AutomorphicArithmetic/AutomorphicVarieties/ModularityArgument/ThresholdArguments/DenseThresholdCases/KeepNumU.lean}{79}{count with the rank-two alternative}.
Both theorems use an additional local comparison witness and give
normalized good-place Frobenius polynomials as well as fixed vectors.
Proposition~\ref{newgal:level} establishes the fixed-vector and weight
conclusions for the stated classical representations. Its rank-two
calculation is over $\mathbb Q$ before restriction to the CM field.

\subsection{The polarized adjoint Selmer group over \texorpdfstring{$\mathbb Q$}{Q}}

Let $K$ be imaginary quadratic, $n\ge2$, $k=\overline{\mathbb Q}_\ell$,
and fix $\iota:k\simeq\mathbb C$. Write $\varepsilon$ for the $\ell$-adic
cyclotomic character, $\delta$ for the quadratic character of $K/\mathbb Q$,
and $c\in G_{\mathbb Q}$ for complex conjugation at a chosen embedding.
Given $s:G_K\to\mathrm{GL}_n(k)$, a multiplier $\mu$ on $G_K$ and
$P\in\mathrm{GL}_n(k)$ with
\[
 s(g)^t P s(cgc^{-1})=\mu(g)P ,
\]
let $A$ be a $G_{\mathbb Q}$-action on $M_n(k)$ which acts by conjugation
by $s(g)$ for $g\in G_K$ and by $X\mapsto-P^{-1}X^tP$ for $c$. In
particular its scalar line affords $\delta$.

\begin{proposition}[The two adjoint vanishing assertions]\label{newgal:adselmer}
Assume $s$ is continuous and absolutely irreducible on
$G_{K(\mu_{\ell^\infty})}$. Suppose one of the following holds, with
attachment meaning equality of normalized Frobenius polynomials at all
but finitely many places for a classical regular algebraic cuspidal
representation $\Pi$ of $\mathrm{GL}_n(\mathbb A_K)$.
\begin{enumerate}
\item $s$ is attached to a conjugate self-dual $\Pi$ and
$\mu=\varepsilon^{1-n}|_{G_K}$.
\item $s=r|_{G_K}$ for a continuous $r:G_{\mathbb Q}\to\mathrm{GL}_n(k)$
and $s$ is attached to $\Pi$. There is a nondegenerate pairing on $r$
with multiplier $\theta=\chi\varepsilon^w$, where $w\in\mathbb Z$ and
$\chi$ factors through $\operatorname{Gal}(\mathbb Q(\mu_M)/\mathbb Q)$
for some $M\ge1$. Here $\mu=\theta|_{G_K}$.
\end{enumerate}
Then $H^1_f(\mathbb Q,A)=0$ for every such $A$.
\end{proposition}
\begin{proof}
In the first case $(\Pi,\delta^n)$ is polarized, since $\delta$ is
trivial on norms from $K$ and $\delta_\infty(-1)=-1$. In the second,
$s^c\simeq s$ by $r(c)$ and $s\simeq s^\vee\otimes\mu$ by the pairing.
Let $\xi$ be the finite order Hecke character of $\mathbb Q$ with
$\ell$-adic realization $\chi$, using geometric reciprocity. Put
$\eta=\xi|\cdot|_{\mathbb A}^{w+n-1}$. Its $\ell$-adic realization is
$\theta\varepsilon^{n-1}$. If the normalized eigenvalues of $\Pi_v$
are $b_i=q_v^{(n-1)/2}\alpha_i$ at a good place, the dual has normalized
eigenvalues $q_v^{n-1}/b_i$. Twisting the dual by $\eta\circ N_{K/\mathbb Q}$
multiplies these by $\iota(\theta(\mathrm{Frob}_v))q_v^{1-n}$.
The result is the multiset for $\Pi^c_v$, by the two isomorphisms above.
Strong multiplicity one gives
$\Pi^c\simeq\Pi^\vee\otimes(\eta\circ N_{K/\mathbb Q}\circ\det)$.
Multiply $\eta$ by $\delta$ if necessary to arrange
$\eta_\infty(-1)=(-1)^n$. The twist on $\mathrm{GL}_n(\mathbb A_K)$
is unchanged, and the Galois multiplier
$\varepsilon^{1-n}\eta_\ell$ restricts to $\mu$ on $G_K$.

In both cases Chebotarev and Brauer--Nesbitt identify $s$ with
$r_{\Pi,\iota}$. Fixing the same $c$, the automorphic polarization gives
a pairing with multiplier $\mu$ and the negative adjoint action above
(Newton--Thorne,
\emph{Adjoint Selmer groups of automorphic Galois representations of unitary
type}, \S1~\cite{newtonThorne}). By Schur's lemma the pairing matrices differ by a scalar,
so they give the same action of $c$. Now apply
Thorne, \emph{On the vanishing
of adjoint Bloch--Kato Selmer groups of irreducible automorphic Galois
representations}, Theorem~6.1~\cite{thorneAdjoint}, with $F=K$ and $F^+=\mathbb Q$.
The hypothesis on the cyclotomic kernel is precisely its irreducibility
hypothesis. This theorem covers $\ell=2$ and imposes no ordinarity or
residual image condition. The earlier Newton--Thorne Theorem~5.2
assumes enormous image on the cyclotomic kernel.

To match the local conditions, a continuous cocycle $z$ and the action
$A$ descend simultaneously to a finite extension $L/\mathbb Q_\ell$:
apply Lemma~\ref{newgal:coeff} to the representation with diagonal
blocks $A(g),1$ and upper right block $z(g)$. Continuity of $A$ follows
from its restriction to the open subgroup $G_K$. At $v\nmid\ell$
the prescribed equation $z(g)=A(g)x-x$ on inertia is the unramified
condition. At $\ell$ the equation is $z(g)\otimes1=gx-x$ in
$A\otimes_{\mathbb Q_\ell}B_{\mathrm{cris}}$, the Bloch--Kato finite
condition. Indeed, for an embedding
$j:\overline{\mathbb Q}\hookrightarrow\overline{\mathbb Q}_\ell$,
the subgroup of compatible pairs of local and global automorphisms
projects isomorphically to $G_{\mathbb Q_\ell}$: restriction to the
algebraic closure of $\mathbb Q$ inside $\overline{\mathbb Q}_\ell$
determines the global component. Other choices conjugate the local
condition. In either equation an $L$-linear retraction $k\to L$
fixing $L$, applied coefficientwise and tensored with the identity
on the period ring when needed, gives a witness over $L$. Outside a
finite set containing $\ell$ and the ramification of $A$, inertia
acts trivially on $A$, so the local equation gives $z|_{I_v}=0$.
The extension representation thus factors through the Galois group
unramified outside this set. Hence $z$ represents a Selmer class with
coefficients in $L$, and vanishing gives a global coboundary.
\end{proof}

In rank three the cyclotomic twist cancels from the adjoint, and its
square supplies the multiplier $\varepsilon^{-2}$. The rank-two block
has the alternating determinant pairing and $w=0$. In each case the
cyclotomic-kernel irreducibility used in Proposition~\ref{newgal:adselmer}
is a separate input to the classical argument. The exported rank-two
proof uses its own cyclotomic irreducibility result, the cuspidal
attachment of the modular block, and additional local witnesses in
\src{src/Development/AutomorphicArithmetic/AutomorphicVarieties/ModularityArgument/ThresholdArguments/DenseThresholdCases/RRRigII.lean}{25}{the second rigidity case}.
The exported rank-four automorphy predicate supplies both a cuspidal
attachment over each imaginary quadratic field on which the restriction
is irreducible and a further comparison witness.
\src{src/Development/AutomorphicArithmetic/AutomorphicVarieties/ModularityArgument/AutomorphyLifting/IsAutCsdA.lean}{16}{Quadratic restriction and its witness}.
At a place above \(\ell\) where the ordinary flag characters
\(\kappa_i\) satisfy \(\kappa_i\ne\kappa_j\varepsilon\) for
all \(i,j\), this witness requires a flag for which
\(\kappa_i^m\ne(\kappa_j\varepsilon)^m\) for every \(m\ge1\).
It also requires, for every regular algebraic cuspidal attachment
\(\Pi\), a single real \(t\) such that at every place away from
\(\ell\) and every place above \(\ell\) satisfying the latter
condition, its cuspidal support consists of characters with
uniformizer values of absolute value strictly between
\(q_v^{t-1/2}\) and \(q_v^{t+1/2}\). Thus the first assertion
involves finite-order ratios too, and the second uses the same
\(t\) across the indicated places.
\src{src/Development/AutomorphicArithmetic/AutomorphicVarieties/ModularityArgument/AutomorphyLifting/WeakOfAtt.lean}{27}{The two requirements on the attachment}
and \src{src/Development/AutomorphicArithmetic/AutomorphicVarieties/ModularityArgument/LeviSmoothLocalization/WinUpDown.lean}{38}{The local support and its absolute-value bounds}.
If $J$ is the alternating similitude form, the pairing
matrix is $P=Jr(c)^{-1}$. Write $X^*=J^{-1}X^tJ$. The map
$X\mapsto (X-X^*)/2$ projects $A$ equivariantly onto the usual
conjugation action on $\mathfrak{sp}(J)$. Indeed conjugation by a
similitude commutes with $*$, and $c$ acts on $A$ by
$X\mapsto-r(c)X^*r(c)^{-1}$. Thus a finite symplectic tangent cocycle
which is a coboundary in $A$ has a primitive in $\mathfrak{sp}(J)$.
This division by two takes place in $k$. On the exported rank-four path,
\src{src/Development/AutomorphicArithmetic/AutomorphicVarieties/ModularityArgument/ThresholdArguments/DenseThresholdCases/ExistsCmFieldTwistH1.lean}{28}{descent from a linked CM tower}
produces an imaginary quadratic field and polarized adjoint vanishing
from an ordinary adjoint lifting theorem. The
\src{src/Development/AutomorphicArithmetic/AutomorphicVarieties/ModularityArgument/ThresholdArguments/DenseThresholdCases/RigSp4.lean}{181}{rank-four rigidity theorem}
combines this with the finite local conditions for the tangent cocycle
and the symplectic projection above. Proposition~\ref{newgal:adselmer}
identifies that vanishing for the classical points satisfying its
hypotheses. The exported lifting theorem concerns points of an ordinary
determinant ring and also uses their local flags, genericity and weight
conditions; establishing the hypotheses for those points is a separate
step.
